


\section{Related Work} \label{sec:related_work}
% General equivariance romero2021group
Group equivariant neural networks have demonstrated their effectiveness in a wide variety of tasks \citep{cohen2016group,cohen2017steerable,weiler2019e2equivariant,rezende2019equivariantflows,romero2021group}.
Recently, various forms and methods to achieve E(3) or SE(3) equivariance have been proposed. \citet{thomas2018tensor, fuchs2020se} utilize the spherical harmonics to compute a basis for the transformations, which allows transformations between higher-order representations. A downside to this method is that the spherical harmonics need to be recomputed which can be expensive. Currently, an extension of this method to arbitrary dimensions is unknown. \citet{finzi2020generalizing} parametrize transformations by mapping kernels on the Lie Algebra. For this method the neural network outputs are in certain situations stochastic, which may be undesirable. \citet{horie2020isometric} proposes a set of isometric invariant and equivairant transformations for Graph Neural Networks.
\citet{kohler2019equivariant_workshop, kohler2020equivariant_icml} propose an $\En$ equivariant network to model 3D point clouds, but the method is only defined for positional data on the nodes without any feature dimensions.

Another related line of research concerns message passing algorithms on molecular data. \cite{gilmer2017neural} presented a message passing setting (or Graph Neural Network) for quantum chemistry, this method is permutation equivariant but not translation or rotation equivariant. \cite{kondor2018covariant} extends the equivariance of GNNs such that each neuron transforms in a specific way under permutations, but this extension only affects its permutation group and not translations or rotations in a geometric space. Further works \cite{schutt2017schnet, schutt2017quantum} build $\En$ invariant message passing networks by inputting the relative distances between points. \citet{klicpera2020directional, klicpera2020fast} in addition to relative distances it includes a modified message passing scheme analogous to Belief Propagation that considers angles and directional information equivariant to rotations. It also uses Bessel functions and spherical harmonics to construct and orthogonal basis. \citet{anderson2019cormorant, miller2020relevance} include $SO(3)$ equivariance in its intermediate layers for modelling the behavior and properties of molecular data. Our method is also framed as a message passing framework but in contrast to these methods it achieves $\En$ equivariance.


\subsection*{Relationship with existing methods}
 In Table \ref{tab:related_work} the EGNN equations are detailed together with some of its closest methods from the literature under the message passing notation from \cite{gilmer2017neural}. This table aims to provide a simple way to compare these different algorithms. It is structured in three main rows that describe i) the edge ii) aggregation and iii) node update operations. The GNN algorithm is the same as the previously introduced in Section \ref{sec:background_equivariance}. Our EGNN algorithm is also equivalent to the description in Section \ref{sec:method_main} but notation has been modified to match the (edge, aggregation, node) format. In all equations $\rmr_{ij}^l = (\rmx_i - \rmx_j)^l$. Notice that except the EGNN, all algorithms have the same aggregation operation and the main differences arise from the edge operation. The algorithm that we call "Radial Field" is the $\En$ equivariant update from \cite{kohler2019equivariant_workshop}. This method is $\En$ equivariant, however its main limitation is that it only operates on $\rmx$ and it doesn't propagate node features $\rmh$ among nodes. In the method $\phi_{rf}$ is modelled as an MLP. Tensor Field Networks (TFN) \cite{thomas2018tensor} instead propagate the node embeddings $\rmh$ but it uses spherical harmonics to compute its learnable weight kernel $ \mathbf{W}^{\ell k}: \mathbb{R}^{3} \rightarrow \mathbb{R}^{(2 \ell+1) \times(2 k+1)}$ which preserves $SE(3)$ equivariance but is expensive to compute an limited to the 3 dimensional space. The SE(3) Transformer \cite{fuchs2020se} (not included in this table), can be interpreted as an extension of TFN with attention. Schnet \cite{schutt2017schnet} can be interpreted as an $\En$ invariant Graph Neural Network where $\phi_{cf}$ receives as input relative distances and outputs a continuous filter convolution that multiplies the neighbor embeddings $\rmh$. Our EGNN differs from these other methods in terms that it performs two different updates in each of the table rows, one related to the embeddings $\rmh$ and another related to the coordinates $\rmx$, these two variables exchange information in the edge operation. In summary the EGNN can retain the flexibility of GNNs while remaining $\En$ equivariant as the Radial Field algorithm and without the need to compute expensive operations (i.e. spherical harmonics).